% ============================================================================
% ANEXOS
% ============================================================================

\chapter*{Anexos}
\addcontentsline{toc}{chapter}{Anexos}

% ============================================================================
% TABLA DE NOMENCLATURA (ATLAS)
% ============================================================================

\section{Nomenclatura y Simbología Matemática}
\label{sec:nomenclatura}

Esta sección presenta la nomenclatura matemática utilizada en el sistema de inferencia difusa y análisis estadístico.

\subsection{Símbolos Matemáticos Generales}

% Tabla multi-página con longtable (permite continuación automática)
\begin{longtable}{@{}p{3cm}p{10cm}@{}}
\caption{Símbolos y Notación Matemática del Sistema Difuso}
\label{tab:nomenclatura} \\
\toprule
\textbf{Símbolo} & \textbf{Descripción} \\
\midrule
\endfirsthead

\multicolumn{2}{c}{\textit{(Continuación de Tabla \ref{tab:nomenclatura})}} \\
\toprule
\textbf{Símbolo} & \textbf{Descripción} \\
\midrule
\endhead

\midrule
\multicolumn{2}{r}{\textit{Continúa en la siguiente página...}} \\
\endfoot

\bottomrule
\endlastfoot
\multicolumn{2}{@{}l}{\textit{\textbf{Conjuntos y Funciones de Membresía}}} \\
\(\tilde{A}\) & Conjunto difuso \\
\(\mu_{\tilde{A}}(x)\) & Función de membresía del conjunto difuso \(\tilde{A}\) \\
\(X\) & Universo de discurso (conjunto clásico de valores posibles) \\
\(X_i\) & Variable de entrada \(i\) (\(i \in \{1,2,3,4\}\)) \\
\(Y\) & Variable de salida (Índice de Sedentarismo) \\
\(T_i\) & Conjunto de términos lingüísticos de variable \(i\) \\
\(\mathbf{P}_i\) & Matriz de percentiles \(3 \times 3\) para variable \(X_i\) \\
\(p_{k,i}\) & Percentil \(k\) de la variable \(X_i\) \\
\(\mathbf{M}^{(j)}\) & Matriz de membresías \(4 \times 3\) para observación \(j\) \\
\\
\multicolumn{2}{@{}l}{\textit{\textbf{Variables Fisiológicas}}} \\
\(\text{Act}_{\text{rel}}\) & Actividad Relativa (min movimiento / hrs monitoreadas) \\
\(\text{Sup}_{\text{cal}}\) & Superávit Calórico Basal (\% del TMB) \\
\(\text{HRV-SDNN}\) & Variabilidad Cardíaca (desviación estándar intervalos RR) \\
\(\Delta_{\text{FC}}\) & Delta Cardíaco (FC caminata - FC reposo) \\
\(\text{TMB}\) & Tasa Metabólica Basal (kcal/día, ec. Mifflin-St Jeor) \\
\(\text{GCA}\) & Gasto Calórico Activo (kcal/día) \\
\(\text{FC}\) & Frecuencia Cardíaca (latidos por minuto) \\
\(\text{IQR}\) & Rango Intercuartílico (Q3 - Q1) \\
\\
\multicolumn{2}{@{}l}{\textit{\textbf{Reglas de Inferencia y Activación}}} \\
\(R_r\) & Regla de inferencia \(r\) (\(r \in \{1,2,3,4,5\}\)) \\
\(\mathcal{R}\) & Base de reglas (conjunto de 5 reglas Mamdani) \\
\(w_r^{(j)}\) & Activación (firing strength) de regla \(R_r\) para observación \(j\) \\
\(\mathbf{w}^{(j)}\) & Vector de activaciones \(\in \mathbb{R}^5\) \\
\(T(a, b)\) & T-norm (operador AND fuzzy): \(T(a,b) = \min(a,b)\) \\
\(y_r\) & Output crisp de la regla \(R_r\) \\
\(\mathbf{y}_{\text{outputs}}\) & Vector de outputs de reglas \([1.0, 0.0, 0.9, 0.5, 0.7]^T\) \\
\\
\multicolumn{2}{@{}l}{\textit{\textbf{Defuzzificación y Clasificación}}} \\
\(y^{(j)}\) & Score de sedentarismo defuzzificado para semana \(j\) \\
\(\hat{y}^{(j)}\) & Clasificación binaria (0=Bajo, 1=Alto sedentarismo) \\
\(\tau\) & Umbral de decisión para binarización \\
\(\|\mathbf{w}\|_1\) & Norma L1 del vector (suma de componentes) \\
\\
\multicolumn{2}{@{}l}{\textit{\textbf{Validación LOOU}}} \\
\(\mathcal{D}\) & Dataset completo (1,337 semanas, 10 usuarios) \\
\(\mathcal{D}_{\text{train}}^{(i)}\) & Conjunto de entrenamiento en fold \(i\) (9 usuarios) \\
\(\mathcal{D}_{\text{test}}^{(i)}\) & Conjunto de prueba en fold \(i\) (1 usuario) \\
\(n_i\) & Número de semanas del usuario \(i\) \\
\(F1^{(i)}\) & F1-Score en fold \(i\) (usuario \(i\) como test) \\
\(\overline{F1}_{\text{LOOU}}\) & F1-Score promedio LOOU (10 folds) \\
\(\sigma_{F1}\) & Desviación estándar de F1-Score entre folds \\
\(CV\) & Coeficiente de variación (\%) \\
\\
\multicolumn{2}{@{}l}{\textit{\textbf{Clustering (Verdad Operativa)}}} \\
\(K\) & Número de clusters (K=2) \\
\(\mathbf{C}_k\) & Centroide del cluster \(k\) (\(k \in \{0,1\}\)) \\
\(c^{(j)}\) & Cluster asignado a semana \(j\) (verdad operativa) \\
\(S\) & Silhouette Score (métrica de calidad clustering) \\
\\
\multicolumn{2}{@{}l}{\textit{\textbf{Métricas de Evaluación}}} \\
\(\text{TP}\) & True Positives (verdaderos positivos) \\
\(\text{TN}\) & True Negatives (verdaderos negativos) \\
\(\text{FP}\) & False Positives (falsos positivos) \\
\(\text{FN}\) & False Negatives (falsos negativos) \\
\(\text{Precision}\) & Precisión: \(\text{TP}/(\text{TP}+\text{FP})\) \\
\(\text{Recall}\) & Sensibilidad: \(\text{TP}/(\text{TP}+\text{FN})\) \\
\(F1\) & F1-Score: \(2 \cdot \text{Prec} \cdot \text{Rec} / (\text{Prec} + \text{Rec})\) \\
\(\text{MCC}\) & Matthews Correlation Coefficient \\
\(\text{Acc}\) & Accuracy (exactitud): \((\text{TP}+\text{TN})/N\) \\
\\
\multicolumn{2}{@{}l}{\textit{\textbf{Operadores y Notación Matemática}}} \\
\(\min(a, b)\) & Mínimo de \(a\) y \(b\) \\
\(\max(a, b)\) & Máximo de \(a\) y \(b\) \\
\(\mathbb{E}[\cdot]\) & Valor esperado (esperanza matemática) \\
\(\mathbb{1}[\cdot]\) & Función indicadora (1 si verdadero, 0 si falso) \\
\(\|\mathbf{v}\|_2\) & Norma euclidiana (L2) del vector \(\mathbf{v}\) \\
\(\|\mathbf{v}\|_1\) & Norma L1 (suma de valores absolutos) \\
\(\mathbb{R}\) & Conjunto de números reales \\
\(\mathbb{R}^+\) & Conjunto de números reales positivos \\
\(\mathbb{R}^{n}\) & Espacio euclidiano n-dimensional \\
\([a, b]\) & Intervalo cerrado de \(a\) a \(b\) \\
\(\arg\min_x f(x)\) & Argumento que minimiza la función \(f(x)\) \\
\(\forall\) & Para todo (cuantificador universal) \\
\(\exists\) & Existe (cuantificador existencial) \\
\(\in\) & Pertenece a (relación de pertenencia) \\
\(\subset\) & Subconjunto de \\
\(\setminus\) & Diferencia de conjuntos \\
\(\land\) & Conjunción lógica (AND) \\
\(\Rightarrow\) & Implicación lógica (THEN) \\
\end{longtable}

% ============================================================================
% FIN DE NOMENCLATURA (ATLAS)
% ============================================================================

% Aquí se incluirán otros anexos de la tesis
% Por ejemplo: consentimiento informado, instrumentos, etc.



